\section{0/1 Knapsack}

\subsection{De bai}

Cho $N$ vat pham, vat thu $i$ co trong luong $w[i]$ va gia tri $v[i]$.
Ban co mot ba lo voi suc chua toi da $W$.
Hay chon mot tap con cac vat pham sao cho tong trong luong khong vuot qua $W$
va tong gia tri la lon nhat.

Rang buoc: moi vat pham chi duoc chon hoac khong chon (0/1), khong duoc chia nho.

\textbf{Input mau:} $N = 4$, $W = 5$, trong luong $w = [2, 3, 4, 5]$, gia tri $v = [3, 4, 5, 6]$.

\subsection{Cong thuc truy hoi}

Goi $dp[i][j]$ la gia tri lon nhat khi xet $i$ vat pham dau tien voi suc chua $j$.
Ta co cong thuc truy hoi:

$$
dp[i][j] = \max\bigl(dp[i-1][j],\; dp[i-1][j - w[i]] + v[i]\bigr)
$$

Truong hop co so: $dp[0][j] = 0$ voi moi $j$ (khong chon vat nao thi gia tri bang $0$).

Dieu kien: chi co the lay vat $i$ khi $j \geq w[i]$. Neu $j < w[i]$ thi $dp[i][j] = dp[i-1][j]$.

\subsection{So do bang DP -- trang thai ban dau}

\begin{diagram}[id="knapsack-init"]
\shape{dp}{DPTable}{rows=5, cols=6, row_labels="i=0..4", col_labels="j=0..5", label="dp[i][j]"}
\apply{dp.cell[0][0]}{value=0}
\apply{dp.cell[0][1]}{value=0}
\apply{dp.cell[0][2]}{value=0}
\apply{dp.cell[0][3]}{value=0}
\apply{dp.cell[0][4]}{value=0}
\apply{dp.cell[0][5]}{value=0}
\recolor{dp.cell[0][0]}{state=done}
\recolor{dp.cell[0][1]}{state=done}
\recolor{dp.cell[0][2]}{state=done}
\recolor{dp.cell[0][3]}{state=done}
\recolor{dp.cell[0][4]}{state=done}
\recolor{dp.cell[0][5]}{state=done}
\end{diagram}

Hang $i = 0$ tuong ung voi viec khong chon vat nao. Moi o $dp[0][j] = 0$.
Cac hang tiep theo se duoc dien bang cach ap dung cong thuc truy hoi,
xet tung vat pham mot tu trai sang phai.

\subsection{Giai thich trang thai DP}

Moi o $dp[i][j]$ dai dien cho bai toan con: \textit{gia tri lon nhat co the dat duoc
khi chi xet $i$ vat pham dau va suc chua la $j$}. Tai moi o, ta co hai lua chon:

\begin{itemize}
\item \textbf{Bo qua} vat $i$: lay gia tri tu hang tren cung cot, $dp[i-1][j]$.
\item \textbf{Lay} vat $i$: lay gia tri $dp[i-1][j - w[i]] + v[i]$, chi khi $j \geq w[i]$.
\end{itemize}

Gia tri lon hon trong hai lua chon se duoc ghi vao $dp[i][j]$.

\subsection{Qua trinh dien bang DP}

\begin{animation}[id="knapsack-fill", label="0/1 Knapsack -- Dien bang DP"]
\shape{items}{Array}{size=4, data=["w=2,v=3","w=3,v=4","w=4,v=5","w=5,v=6"], labels="1..4", label="Vat pham"}
\shape{dp}{DPTable}{rows=5, cols=6, row_labels="i=0..4", col_labels="j=0..5", label="dp[i][j]"}

\step
\apply{dp.cell[0][0]}{value=0}
\apply{dp.cell[0][1]}{value=0}
\apply{dp.cell[0][2]}{value=0}
\apply{dp.cell[0][3]}{value=0}
\apply{dp.cell[0][4]}{value=0}
\apply{dp.cell[0][5]}{value=0}
\recolor{dp.cell[0][0]}{state=done}
\recolor{dp.cell[0][1]}{state=done}
\recolor{dp.cell[0][2]}{state=done}
\recolor{dp.cell[0][3]}{state=done}
\recolor{dp.cell[0][4]}{state=done}
\recolor{dp.cell[0][5]}{state=done}
\narrate{Co so: $dp[0][j] = 0$ voi moi $j$. Khong co vat pham nao duoc xet nen gia tri luon bang $0$.}

\step
\recolor{items.cell[0]}{state=current}
\recolor{dp.cell[1][0]}{state=current}
\recolor{dp.cell[1][1]}{state=current}
\annotate{dp.cell[1][0]}{label="skip", arrow_from="dp.cell[0][0]", color=info}
\annotate{dp.cell[1][1]}{label="skip", arrow_from="dp.cell[0][1]", color=info}
\apply{dp.cell[1][0]}{value=0}
\apply{dp.cell[1][1]}{value=0}
\narrate{Vat $1$ ($w=2, v=3$). Voi $j=0$ va $j=1$: suc chua khong du ($j < 2$), chi co the bo qua. $dp[1][0] = 0$, $dp[1][1] = 0$.}

\step
\recolor{dp.cell[1][0]}{state=done}
\recolor{dp.cell[1][1]}{state=done}
\recolor{dp.cell[1][2]}{state=current}
\annotate{dp.cell[1][2]}{label="skip: 0", arrow_from="dp.cell[0][2]", color=muted}
\annotate{dp.cell[1][2]}{label="take: 0+3=3", arrow_from="dp.cell[0][0]", color=good}
\narrate{$j=2$: bo qua cho $dp[0][2]=0$. Lay vat $1$: $dp[0][2-2]+3 = 0+3 = 3$. $\max(0, 3) = 3$. \textbf{Lay!}}

\step
\apply{dp.cell[1][2]}{value=3}
\recolor{dp.cell[1][2]}{state=done}
\recolor{dp.cell[1][3]}{state=current}
\annotate{dp.cell[1][3]}{label="skip: 0", arrow_from="dp.cell[0][3]", color=muted}
\annotate{dp.cell[1][3]}{label="take: 0+3=3", arrow_from="dp.cell[0][1]", color=good}
\narrate{$j=3$: bo qua cho $0$, lay cho $dp[0][1]+3 = 3$. $\max(0,3) = 3$. \textbf{Lay!}}

\step
\apply{dp.cell[1][3]}{value=3}
\recolor{dp.cell[1][3]}{state=done}
\recolor{dp.cell[1][4]}{state=current}
\recolor{dp.cell[1][5]}{state=current}
\annotate{dp.cell[1][4]}{label="take: 3", arrow_from="dp.cell[0][2]", color=good}
\annotate{dp.cell[1][5]}{label="take: 3", arrow_from="dp.cell[0][3]", color=good}
\apply{dp.cell[1][4]}{value=3}
\apply{dp.cell[1][5]}{value=3}
\narrate{$j=4$: $\max(0, dp[0][2]+3) = 3$. $j=5$: $\max(0, dp[0][3]+3) = 3$. Hang $i=1$ hoan thanh. Vat $1$ duoc lay khi $j \geq 2$.}

\step
\recolor{dp.cell[1][4]}{state=done}
\recolor{dp.cell[1][5]}{state=done}
\recolor{items.cell[0]}{state=done}
\recolor{items.cell[1]}{state=current}
\recolor{dp.cell[2][0]}{state=current}
\recolor{dp.cell[2][1]}{state=current}
\recolor{dp.cell[2][2]}{state=current}
\apply{dp.cell[2][0]}{value=0}
\apply{dp.cell[2][1]}{value=0}
\apply{dp.cell[2][2]}{value=3}
\annotate{dp.cell[2][2]}{label="skip: 3", arrow_from="dp.cell[1][2]", color=good}
\narrate{Vat $2$ ($w=3, v=4$). $j=0,1,2$: suc chua khong du de lay vat $2$. Bo qua, sao chep tu hang tren: $dp[2][0..2] = [0, 0, 3]$.}

\step
\recolor{dp.cell[2][0]}{state=done}
\recolor{dp.cell[2][1]}{state=done}
\recolor{dp.cell[2][2]}{state=done}
\recolor{dp.cell[2][3]}{state=current}
\annotate{dp.cell[2][3]}{label="skip: 3", arrow_from="dp.cell[1][3]", color=muted}
\annotate{dp.cell[2][3]}{label="take: 0+4=4", arrow_from="dp.cell[1][0]", color=good}
\narrate{$j=3$: bo qua cho $3$. Lay vat $2$: $dp[1][0]+4 = 0+4 = 4$. $\max(3, 4) = 4$. \textbf{Lay!}}

\step
\apply{dp.cell[2][3]}{value=4}
\recolor{dp.cell[2][3]}{state=done}
\recolor{dp.cell[2][4]}{state=current}
\annotate{dp.cell[2][4]}{label="skip: 3", arrow_from="dp.cell[1][4]", color=muted}
\annotate{dp.cell[2][4]}{label="take: 0+4=4", arrow_from="dp.cell[1][1]", color=good}
\narrate{$j=4$: bo qua cho $3$. Lay vat $2$: $dp[1][1]+4 = 0+4 = 4$. $\max(3, 4) = 4$. \textbf{Lay!}}

\step
\apply{dp.cell[2][4]}{value=4}
\recolor{dp.cell[2][4]}{state=done}
\recolor{dp.cell[2][5]}{state=current}
\annotate{dp.cell[2][5]}{label="skip: 3", arrow_from="dp.cell[1][5]", color=muted}
\annotate{dp.cell[2][5]}{label="take: 3+4=7", arrow_from="dp.cell[1][2]", color=good}
\narrate{$j=5$: bo qua cho $3$. Lay vat $2$: $dp[1][2]+4 = 3+4 = 7$. $\max(3, 7) = 7$. \textbf{Lay!} Ca hai vat $1$ va $2$ deu duoc chon.}

\step
\apply{dp.cell[2][5]}{value=7}
\recolor{dp.cell[2][5]}{state=done}
\recolor{items.cell[1]}{state=done}
\recolor{items.cell[2]}{state=current}
\apply{dp.cell[3][0]}{value=0}
\apply{dp.cell[3][1]}{value=0}
\apply{dp.cell[3][2]}{value=3}
\apply{dp.cell[3][3]}{value=4}
\recolor{dp.cell[3][0]}{state=done}
\recolor{dp.cell[3][1]}{state=done}
\recolor{dp.cell[3][2]}{state=done}
\recolor{dp.cell[3][3]}{state=done}
\recolor{dp.cell[3][4]}{state=current}
\annotate{dp.cell[3][4]}{label="skip: 4", arrow_from="dp.cell[2][4]", color=good}
\annotate{dp.cell[3][4]}{label="take: 0+5=5", arrow_from="dp.cell[2][0]", color=muted}
\narrate{Vat $3$ ($w=4, v=5$). $j=0..3$: khong du suc chua hoac bo qua tot hon, sao chep tu hang tren. $j=4$: bo qua cho $4$, lay cho $dp[2][0]+5=5$. $\max(4,5) = 5$. \textbf{Lay!}}

\step
\apply{dp.cell[3][4]}{value=5}
\recolor{dp.cell[3][4]}{state=done}
\recolor{dp.cell[3][5]}{state=current}
\annotate{dp.cell[3][5]}{label="skip: 7", arrow_from="dp.cell[2][5]", color=good}
\annotate{dp.cell[3][5]}{label="take: 0+5=5", arrow_from="dp.cell[2][1]", color=muted}
\narrate{$j=5$: bo qua cho $7$, lay cho $dp[2][1]+5 = 0+5 = 5$. $\max(7, 5) = 7$. \textbf{Bo qua!} Giu nguyen to hop vat $1+2$ tot hon.}

\step
\apply{dp.cell[3][5]}{value=7}
\recolor{dp.cell[3][5]}{state=done}
\recolor{items.cell[2]}{state=done}
\recolor{items.cell[3]}{state=current}
\apply{dp.cell[4][0]}{value=0}
\apply{dp.cell[4][1]}{value=0}
\apply{dp.cell[4][2]}{value=3}
\apply{dp.cell[4][3]}{value=4}
\apply{dp.cell[4][4]}{value=5}
\recolor{dp.cell[4][0]}{state=done}
\recolor{dp.cell[4][1]}{state=done}
\recolor{dp.cell[4][2]}{state=done}
\recolor{dp.cell[4][3]}{state=done}
\recolor{dp.cell[4][4]}{state=done}
\recolor{dp.cell[4][5]}{state=current}
\annotate{dp.cell[4][5]}{label="skip: 7", arrow_from="dp.cell[3][5]", color=good}
\annotate{dp.cell[4][5]}{label="take: 0+6=6", arrow_from="dp.cell[3][0]", color=muted}
\narrate{Vat $4$ ($w=5, v=6$). $j=0..4$: sao chep tu hang tren (khong du suc chua hoac bo qua tot hon). $j=5$: bo qua cho $7$, lay cho $dp[3][0]+6 = 6$. $\max(7, 6) = 7$. \textbf{Bo qua!}}

\step
\apply{dp.cell[4][5]}{value=7}
\recolor{dp.cell[4][5]}{state=good}
\recolor{items.cell[3]}{state=done}
\recolor{dp.cell[2][5]}{state=path}
\recolor{dp.cell[1][2]}{state=path}
\recolor{dp.cell[0][0]}{state=path}
\reannotate{dp.cell[2][5]}{color=path, arrow_from="dp.cell[1][2]"}
\reannotate{dp.cell[1][2]}{color=path, arrow_from="dp.cell[0][0]"}
\narrate{$dp[4][5] = 7$. Dap an: gia tri lon nhat la $7$. Truy vet: $dp[4][5] = dp[3][5] = dp[2][5]$ (lay vat $2$ tai day) $\to dp[1][2]$ (lay vat $1$ tai day) $\to dp[0][0]$. Chon vat $1$ va $2$, tong trong $2+3=5$, tong gia tri $3+4=7$.}
\end{animation}

\subsection{Truy vet loi giai}

Sau khi dien xong bang DP, dap an nam o $dp[N][W]$.
De biet nhung vat nao duoc chon, ta truy vet tu $dp[N][W]$ nguoc ve $dp[0][0]$:

\begin{enumerate}
\item Bat dau tai $i = N$, $j = W$.
\item Neu $dp[i][j] \neq dp[i-1][j]$, thi vat $i$ duoc chon. Giam $j$ di $w[i]$, giam $i$ di $1$.
\item Neu $dp[i][j] = dp[i-1][j]$, thi vat $i$ khong duoc chon. Chi giam $i$ di $1$.
\item Lap lai cho den khi $i = 0$.
\end{enumerate}

Trong vi du tren: $dp[4][5] = 7 = dp[3][5]$ (khong lay vat $4$),
$dp[3][5] = 7 = dp[2][5]$ (khong lay vat $3$),
$dp[2][5] = 7 \neq dp[1][5] = 3$ (lay vat $2$, $j$ giam tu $5$ ve $2$),
$dp[1][2] = 3 \neq dp[0][2] = 0$ (lay vat $1$, $j$ giam tu $2$ ve $0$).
Vay tap vat duoc chon la $\{1, 2\}$ voi tong gia tri $3 + 4 = 7$.

\subsection{Cai dat}

\begin{lstlisting}[language=C++]
#include <bits/stdc++.h>
using namespace std;

int main() {
    int n, W;
    cin >> n >> W;
    vector<int> w(n + 1), v(n + 1);
    for (int i = 1; i <= n; i++)
        cin >> w[i] >> v[i];

    // dp[i][j] = max value using items 1..i with capacity j
    vector<vector<int>> dp(n + 1, vector<int>(W + 1, 0));

    for (int i = 1; i <= n; i++) {
        for (int j = 0; j <= W; j++) {
            dp[i][j] = dp[i - 1][j]; // skip item i
            if (j >= w[i])
                dp[i][j] = max(dp[i][j], dp[i - 1][j - w[i]] + v[i]);
        }
    }

    cout << dp[n][W] << "\n";

    // Backtrack to find chosen items
    vector<int> chosen;
    for (int i = n, j = W; i >= 1; i--) {
        if (dp[i][j] != dp[i - 1][j]) {
            chosen.push_back(i);
            j -= w[i];
        }
    }

    reverse(chosen.begin(), chosen.end());
    for (int x : chosen)
        cout << x << " ";
    cout << "\n";

    return 0;
}
\end{lstlisting}

Do phuc tap: $O(N \cdot W)$ thoi gian va bo nho.
Co the toi uu bo nho xuong $O(W)$ bang cach chi giu mot hang DP va duyet $j$ giam dan.
